% F1 -- What a single prediction is. Deliberately the simplest figure in the
% report: three columns A, B, C, and two rows, one per kind of moment.
\begin{figure}[!htbp]
\centering
\fitwidth{%
\begin{tikzpicture}[font=\small,
  ev/.style={draw=clrule,fill=clsoft,rounded corners=2pt,inner sep=3pt,
             font=\scriptsize\ttfamily,minimum width=1.35cm,minimum height=0.55cm},
  tgt/.style={draw=clmodel,fill=white,thick,rounded corners=2pt,inner sep=3pt,
              minimum width=1.35cm,minimum height=0.55cm},
  hr/.style={font=\scriptsize,text=clrule},
  hdr/.style={font=\footnotesize\bfseries,text=clmodel,anchor=west},
  cs/.style={font=\footnotesize\bfseries,anchor=east},
  ans/.style={font=\scriptsize,anchor=west},
  note/.style={font=\scriptsize,text=clrule,anchor=west},
]
% ================= column headers =================
\node[hdr] at (-0.70,2.05) {A. What the model is given};
\node[hdr] at (5.55,2.05)  {B. The moment};
\node[hdr] at (8.30,2.05)  {C. The two answers, compared};
\draw[clrule] (-0.80,1.80) -- (14.2,1.80);

% ================= row 1: the person has not moved =================
\node[cs] at (-1.00,0.85) {1};
\foreach \i/\h/\c in {0/07h/BKVCHE1, 1/09h/BKVCHE1, 2/11h/BKVCHE1, 3/12h/BKVCHE1}{
  \node[hr] at (\i*1.50,1.42) {\h};
  \node[ev] at (\i*1.50,0.95) {\c};
}
\node[hr]  at (6.15,1.42) {13h};
\node[tgt] at (6.15,0.95) {\textbf{?}};
\node[ans] at (8.30,1.32) {really visited: \cid{BKVCHE1}};
\node[ans] at (8.30,0.97) {the rule says \cid{BKVCHE1}: \textcolor{clgain}{\textbf{right}}};
\node[ans] at (8.30,0.62) {my model says \cid{BKVCHE1}: \textcolor{clgain}{\textbf{right}}};
\node[note] at (-0.70,0.20)
  {\textbf{The person has not moved.} Both answers are free, and neither is
   about the model. This is the ordinary case.};

\draw[clrule,dotted] (-0.80,-0.18) -- (14.2,-0.18);

% ================= row 2: the person has gone somewhere =================
\node[cs] at (-1.00,-1.15) {2};
\foreach \i/\h/\c in {0/07h/BKVCHE1, 1/09h/BKVCHE1, 2/11h/BKVCHE1, 3/12h/BKVCHE1}{
  \node[hr] at (\i*1.50,-0.58) {\h};
  \node[ev] at (\i*1.50,-1.05) {\c};
}
\node[hr]  at (6.15,-0.58) {18h};
\node[tgt] at (6.15,-1.05) {\textbf{?}};
\node[ans] at (8.30,-0.68) {really visited: \cid{BKVPAN1}};
\node[ans] at (8.30,-1.03) {the rule says \cid{BKVCHE1}: \textcolor{clloss}{\textbf{wrong}}};
\node[ans] at (8.30,-1.38) {my model says \cid{BKVPAN1}: \textcolor{clgain}{\textbf{right}}};
\node[note] at (-0.70,-1.80)
  {\textbf{The person has gone somewhere.} The rule must fail. This is the only
   kind of moment on which I can win anything.};
\end{tikzpicture}}
\caption{What one prediction is}
\label{fig:worked-example}
\figtext{%
\figpara{How to read it.}{The two rows are two different moments of the same
person's day, and the three columns are always the same three things. Column A
is the past the model is allowed to read: four moments, each written as an hour
and the name of the antenna the phone was on. Column B is the moment I want an
answer for, with the hour given and the antenna hidden. Column C compares the
answer of the one-line rule with the answer of my model against the antenna the
person really visited.}

\figpara{Row 1: nothing happens.}{The person is still on the same antenna at
13h. Answering ``the same as before'' is correct, so the one-line rule scores a
point without knowing anything, and my model scores the same point. Moments like
this one are the majority, which is why a rule with no model in it is right
about seven times out of ten.}

\figpara{Row 2: the person leaves.}{At 18h they are on a different antenna. The
rule now has to be wrong, because being wrong is all it can do when someone
moves. This is the only kind of moment on which a model can gain anything at
all, and every gain reported in this document was won on moments of this second
kind. The hour is what makes the second row answerable: 18h is when this
particular person leaves.}

\figpara{One caution.}{The two rows are drawn to explain the task, not measured
from a run. The antenna names are real names from the study area; how often each
of the two cases actually occurs is measured in Chapter~\ref{ch:data}.}}
\end{figure}
